\scanpage{169}
\mspage{81}
Il me semble donc plausible que ces topos, avec leurs morphismes canoniques entre eux qu'il faudrait expliciter, \unsure{suffiraient|suffiront} à exprimer tous les topos \inserted{canoniques} associés à la stratification, envisagés pp 74' \unsure{ff.}. \deleted{Prenons} Revenons au cas-type, avec $n=2$, i.e. $\mathbb C^2$, \deleted{\unsure{donc les}}\inserted{il y a} alors $4\times4=16$ \deleted{strates} topos \inserted{élémentaires} qui s'agencent d'une certaine façon :
\margin{$*$}
\diagram{\[
(277)\quad\left\{\begin{array}{l}
\left\{\begin{array}{ll}
D_0^*=\mathbb C^*\times\mathbb C^*&(D_0=X=\mathbb C\times\mathbb C)\\
\left.\begin{array}{l}
\Delta_1^*=\mathbb C^*\times\mathbb P\\
\Delta_2^{\unsure{*|'}}=\mathbb P\times\mathbb C^*
\end{array}\right\}&
\begin{array}{l}
D_1^*=\Delta_1^*\amalg\Delta_2^*\quad(D_1=\Delta_1\amalg\Delta_2)\\
\text{où }\Delta_1=\mathbb C\times\mathbb P\\
\Delta_2=\mathbb P\times\mathbb C
\end{array}\\
\mathbb P=\mathbb P\times\mathbb P&\}\ D_2^*=D_2=\mathbb P
\end{array}\right.\\[2em]
\left\{\begin{array}{ll}
\unsure{\mathcal V}_{\Delta_1^*,X}=\mathbb C^*\times\mathbb D&
\unsure{\mathcal V}^*_{\Delta_1^*,X}=\mathbb C^*\times\mathbb D^*\\
\unsure{\mathcal V}_{\Delta_2^*,X}=\mathbb D\times\mathbb C^*&
\unsure{\mathcal V}^*_{\unsure{\Delta_2^*|\Delta_1^*},X}=\mathbb D^*\times\mathbb C^*
\end{array}\right.\\[2em]
\left\{\begin{array}{ll}
\unsure{\mathcal V}_{\mathbb P,\Delta_1}=\mathbb D\times\mathbb P&
\unsure{\mathcal V}^*_{\mathbb P,\Delta_1}=\mathbb D^*\times\mathbb P\\
\unsure{\mathcal V}^{\unsure{?}}_{\mathbb P,\Delta_2}=\mathbb P\times\mathbb D&
\unsure{\mathcal V}^*_{\mathbb P,\Delta_2}=\mathbb P\times\mathbb D^*
\end{array}\right.\\[2em]
\left\{\begin{array}{ll}
\unsure{\mathcal V}_{\mathbb P,X}=\mathbb D\times\mathbb D,&\\
\unsure{\mathcal V}^{\Delta_1}_{\mathbb P,X}=\mathbb D\times\mathbb D^*,&
\unsure{\mathcal V}^{\Delta_2}_{\mathbb P,X}=\mathbb D^*\times\mathbb D\\
\unsure{\mathcal V}^*_{\mathbb P,X}=\mathbb D^*\times\mathbb D^*&
\end{array}\right.
\end{array}\right.
\]}
\margin{Ces écritures sont bien entendu purement heuristiques et symboliques -- pour les remplacer par des formules correctes$*$, il faudrait pas mal de \unsure{soins} supplémentaires !}
\margin{168}
